%=======================================================================
%  LAST-MINUTE REVISION
%  Master-only. Every count here is the one the chapter itself prints, and
%  each of those was checked against the Detailed PYQ by src\check.py.
%  Nothing in this section is new material: every line points back into a
%  chapter. House style of this subject: no em dashes.
%=======================================================================
\cleardoublepage
\section{Last-Minute Revision}

{\color{sub}\footnotesize Built from the frequency analysis of all 24 papers (2069--2082 BS),
EX 716 (BEI) and EX 752 (BEX), plus EG 785 for 69 Bh. Nothing here is new: every line points
back into a chapter.}

\hr

\T{R.1 Must Know \mk{(every TOP-tier topic: these carry the paper)}}

{\color{sub}\footnotesize A topic is TOP tier when \textbf{7 or more} of the 24 papers ask it.
There are 21 of them and every chapter has at least one, except Chapter 8, whose marks are
spread thin over many small instruments instead.}

\vspace{2pt}
\begin{tabularx}{\textwidth}{@{}L{0.7cm}L{5.5cm}L{1.3cm}X@{}}
\toprule
\hrow \thead{Ch} & \thead{Topic} & \thead{Papers} & \thead{The one thing to be able to do} \\
6 & \textbf{Stability and maximum gain of a transistor} & \textbf{22/24} & $\Delta$, $K$, then $G_{T,max}$; if $K<1$ say \textbf{MSG}, do not quote MAG \\
6 & \textbf{Insertion loss method; Butterworth Vs Chebyshev prototype} & \textbf{20/24} & $P_{LR}$, order from the attenuation curve, $g$-values, scale and transform \\
1 & \textbf{Microwave Vs conventional low-frequency circuits} & \textbf{18/24} & The $d/\lambda$ criterion first, then the comparison table \\
7 & \textbf{How microwaves are hazardous; HERP, HERO, HERF} & \textbf{17/24} & Non-ionizing $\rightarrow$ dielectric heating $\rightarrow$ no internal heat sensors $\rightarrow$ poorly cooled organs \\
2 & \textbf{Double-stub matching, all steps} & \textbf{17/24} & Chart construction with the spacing circle; two lengths and the forbidden region \\
4 & \textbf{Rectangular waveguide TE and TM field equations} & 14/24 & Six steps, then the characteristic-equation table \\
6 & \textbf{Microstrip implementation of the filter} & 13/24 & High-$Z$ line $=$ series $L$, low-$Z$ line $=$ shunt $C$; draw the named section \\
3 & \textbf{Why S-parameters, and S Vs h-parameters} & 12/24 & A matched load exists at every frequency; a short and an open do not \\
7 & \textbf{Radiation fields (zones) of an antenna} & 11/24 & $0.62\sqrt{D^3/\lambda}$ and $2D^2/\lambda$, and the $D>\lambda$ validity check \\
1 & \textbf{Real-world microwave systems and protocols} & 11/24 & Three or five systems with band, modulation, data rate and standard \\
6 & Amplifier design flowchart; stability analysis & 10/24 & Boxes joined by arrows: $\Delta$, $K$, circles, $\Gamma_S$, $\Gamma_L$, networks \\
2 & S-matrix of the network before and after matching & 10/24 & After matching $S_{11}=S_{22}=0$; before, $S_{11}=\Gamma_{in}$ \\
1 & Why S-parameters instead of Z, Y, h & 10/24 & Same answer as Ch 3's; the papers file it under both \\
5 & \textbf{Cavity magnetron and cross-field effect} & 9/24 & $\phi = 2\pi n/N$, Hull cut-off, the $\pi$-mode \\
5 & \textbf{Bunching, velocity and density modulation} & 9/24 & Two lines for 2 marks; $X=1.841$ and $2I_0J_1(X)$ for 10 \\
4 & \textbf{Magic tee and its S-matrix} & 9/24 & Sign pattern down a column, never the port number \\
4 & Dominant and degenerate modes; which cannot exist & 9/24 & $f_c$ ranking; TM$_{m0}$ and TM$_{0n}$ do not exist \\
7 & Safety practices and standards & 7/24 & Name the body, the range and one number each \\
5 & Two-cavity and multicavity klystron & 7/24 & Velocity modulate, drift, bunch, decelerate at the catcher \\
3 & Derive S-parameters for a two-port & 7/24 & Five steps ending at $S_{11}=(Z-Z_0)/(Z+Z_0)$ \\
1 & IEEE band classification and microwave features & 7/24 & The band table, then advantages and disadvantages \\
\bottomrule
\end{tabularx}

\vspace{2pt}
{\footnotesize\color{sub} \mk{Two of these rows are one question.} ``Why S-parameters''
is tagged into Chapter 1 by ten papers and into Chapter 3 by twelve, because the papers
ask it as an opener to both a behaviour question and a network question. Learn it once.}

\par\penalty0\vspace{4pt}

%=======================================================================
\T{R.2 Must Memorize \mk{(what cannot be derived in the hall)}}

\lead{Ch 1: introduction}
\begin{itemize}
  \item \textbf{Microwaves} 300 MHz to 300 GHz, $\lambda$ 1 m down to 1 mm. IEEE letters:
    L 1--2, S 2--4, C 4--8, X 8--12, Ku 12--18, K 18--27, Ka 27--40 GHz.
  \item \textbf{The criterion}: lumped analysis holds while $d/\lambda \ll 1$ (below about
    $0.01$). Above it the line is distributed and $V$, $I$ depend on position.
  \item Parasitics that kill a lumped part: a resistor becomes $R$ with lead $L$ and shunt
    $C$; an inductor self-resonates; a capacitor turns inductive above its SRF.
\end{itemize}

\lead{Ch 2: transmission lines and the Smith chart}
\begin{itemize}
  \item $Z_{in} = Z_0\dfrac{Z_L + jZ_0\tan\beta d}{Z_0 + jZ_L\tan\beta d}$,
    $\Gamma = \dfrac{Z_L-Z_0}{Z_L+Z_0}$, $S = \dfrac{1+|\Gamma|}{1-|\Gamma|}$.
  \item \textbf{Special lengths}: short $\rightarrow jZ_0\tan\beta d$; open $\rightarrow
    -jZ_0\cot\beta d$; $\lambda/2 \rightarrow Z_L$; $\lambda/4 \rightarrow Z_0^2/Z_L$.
  \item \textbf{Chart conventions}, both easy to reverse: the WTG scale reads $0.000$ at
    $\Gamma=-1$ and increases \textbf{clockwise}; on an \textbf{admittance} chart the
    terminations swap ends, open at $\Gamma=-1$, short at $\Gamma=+1$.
  \item \textbf{Double stub}: $B = \dfrac{1 \pm \sqrt{g(1+t^2) - g^2t^2}}{t}$, $t=\tan\beta d$,
    a leading \textbf{1}, not $g$. Forbidden region $g > 1/\sin^2\beta d$.
\end{itemize}

\lead{Ch 3: S-parameters}
\begin{itemize}
  \item $b_i = \sum_j S_{ij}a_j$; $S_{ij} = b_i/a_j$ with every other port
    \textbf{matched}, not shorted or open.
  \item \textbf{Four tests}: $S_{ii}=0$ matched; $[S]=[S]^T$ reciprocal;
    $[S]^\dagger[S]=[I]$ lossless (columns unit, distinct columns orthogonal); reference
    plane shift multiplies $S_{ij}$ by $e^{-j(\theta_i+\theta_j)}$.
  \item \textbf{No 3-port} can be lossless, reciprocal and matched at once. Give up one:
    lossy gives the resistive divider, non-reciprocal the circulator, unmatched the tee.
  \item \textbf{Magic tee} $\dfrac{1}{\sqrt2}$ with the E-arm column carrying $+1, -1$ and
    the H-arm column $+1, +1$. Read the \textbf{sign pattern}, not the port numbers.
\end{itemize}

\lead{Ch 4: waveguides and components}
\begin{itemize}
  \item $f_c = \dfrac{1}{2\sqrt{\mu\epsilon}}\sqrt{(m/a)^2+(n/b)^2}$,
    $\lambda_c = 2a$ for TE$_{10}$; $\lambda_g = \dfrac{\lambda}{\sqrt{1-(f_c/f)^2}}$;
    $\dfrac{1}{\lambda^2} = \dfrac{1}{\lambda_g^2} + \dfrac{1}{\lambda_c^2}$.
  \item $v_p v_g = v^2$, $v_p > v > v_g$. Below $f_c$ the mode is \textbf{evanescent}: a
    guide is a high-pass filter.
  \item $Z_{TE} = \eta/\sqrt{1-(f_c/f)^2}$, $Z_{TM} = \eta\sqrt{1-(f_c/f)^2}$,
    $\eta_0 = 120\pi \approx 377\,\Omega$.
  \item \textbf{TM$_{m0}$ and TM$_{0n}$ do not exist}; TE$_{10}$ is dominant when
    $a > b$. Circular guide: TE$_{11}$ dominant ($p'_{11}=1.841$), TM$_{01}$ next
    ($p_{01}=2.405$).
  \item Directional coupler (input at 1, coupled 4, isolated 3): coupling
    $C = 10\log P_1/P_4$, directivity $D = 10\log P_4/P_3$, isolation
    $I = 10\log P_1/P_3$, and $\boxed{I = C + D}$.
\end{itemize}

\lead{Ch 5: microwave generators}
\begin{itemize}
  \item $v_0 = \sqrt{2eV_0/m} = 0.593\times10^6\sqrt{V_0}$ m/s.
  \item Bunching parameter $X = \dfrac{\beta_i V_1}{2V_0}\theta_0$; catcher fundamental
    $2I_0J_1(X)$, maximum at $X = \mathbf{1.841}$ where it is $1.164I_0$.
  \item \textbf{Reflex klystron} oscillates at $n + 3/4$ cycles in the repeller space;
    best mode efficiency about 22.7\,\%.
  \item \textbf{Magnetron}: adjacent-cavity phase $\phi = 2\pi n/N$, so $\pi$-mode is
    $n = N/2$; Hull cut-off $B_c = \dfrac{\sqrt{8mV_0/e}}{b(1-a^2/b^2)}$.
\end{itemize}

\lead{Ch 6: RF design practices}
\begin{itemize}
  \item $\Delta = S_{11}S_{22} - S_{12}S_{21}$;
    $K = \dfrac{1-|S_{11}|^2-|S_{22}|^2+|\Delta|^2}{2|S_{12}S_{21}|}$. Unconditionally
    stable needs \textbf{$K>1$ and $|\Delta|<1$}, both.
  \item $G_{T,max} = \dfrac{|S_{21}|}{|S_{12}|}\left(K - \sqrt{K^2-1}\right)$, valid
    \textbf{only} when $K>1$. If $K<1$ quote $\mathrm{MSG} = |S_{21}|/|S_{12}|$.
  \item Unilateral: $G_{TU,max} = \dfrac{|S_{21}|^2}{(1-|S_{11}|^2)(1-|S_{22}|^2)}$.
  \item $P_{LR} = 1/(1-|\Gamma|^2)$. Butterworth $g_k = 2\sin\frac{(2k-1)\pi}{2N}$;
    Chebyshev from the ripple. Order from the attenuation curve, never guessed.
  \item Microstrip: a \textbf{short high-impedance} line is a series $L$, a \textbf{short
    low-impedance} line a shunt $C$. That one sentence builds every filter layout.
\end{itemize}

\lead{Ch 7: antennas, propagation and hazards}
\begin{itemize}
  \item Zones: $R_1 = 0.62\sqrt{D^3/\lambda}$, $R_2 = 2D^2/\lambda$, \textbf{both needing
    $D > \lambda$}. For a small antenna use $\lambda/2\pi$ and $2\lambda$ instead.
  \item $S = \dfrac{P G}{4\pi R^2}$, $S = E^2/377$ \textbf{in the far field only}.
  \item $\mathrm{SAR} = \dfrac{\sigma E^2}{\rho}$ W/kg. Whole-body \textbf{4 W/kg} warms the
    body about $1^\circ$C: divide by 10 for workers (0.4), by 50 for the public (0.08).
  \item ICNIRP 1998 public limit, 400 MHz to 2 GHz: $f/200$ W/m$^2$ with $f$ in MHz.
\end{itemize}

\lead{Ch 8: measurements}
\begin{itemize}
  \item \textbf{Bands}: low $<$ 10 mW, medium 10 mW to 10 W, high $>$ 10 W. $P_{av} =
    P_{peak}\times$ duty, duty $=\tau f_r$.
  \item Static calorimeter $P = \dfrac{4.187\,mC_pT}{t}$; circulating
    $P = 4.187\,vdC_p(T_2-T_1)$. \textbf{Same law}, with $m/t$ replaced by $v\,d$.
  \item Double bridge $P_{av} = \dfrac{E_1^2 - E_2^2}{4R}$.
  \item Double minimum, for $S$ above about 10:
    $S = \sqrt{1 + \dfrac{1}{\sin^2(\pi\Delta x/\lambda_g)}} \approx
    \dfrac{\lambda_g}{\pi\Delta x}$.
\end{itemize}

\par\penalty0\vspace{4pt}

%=======================================================================
\T{R.3 Must Practice \mk{(the 56 solved problems, ranked by how many papers set them)}}

{\color{sub}\footnotesize All of these are worked in the companion volume,
\textbf{RF and Microwave Engineering: Complete Numerical Problems and Solutions}. Chapters
1, 5 and 8 set no numerical in 24 papers.}

\vspace{2pt}
\begin{tabularx}{\textwidth}{@{}L{0.7cm}L{5.0cm}L{1.3cm}L{1.5cm}X@{}}
\toprule
\hrow \thead{Ch} & \thead{Family} & \thead{Papers} & \thead{Problems} & \thead{Where it goes wrong} \\
6 & \textbf{Transistor stability and maximum gain} & \textbf{22/24} & 1 to 4 & Quoting MAG when $K<1$. Three of the 13 sets are not unconditionally stable \\
2 & \textbf{Double-stub shunt tuner} & \textbf{17/24} & 8 to 24 & The admittance chart's terminations are mirrored; a wrong end costs $\lambda/4$ \\
6 & Filter design and microstrip layout & 13/24 & 10 & Choosing $N$ by habit rather than from the attenuation curve \\
4 & \textbf{Rectangular guide parameter sets} & 14/24 & 1 to 3 & Using $\lambda$ where $\lambda_g$ belongs; forgetting the cut-off check \\
4 & Magic tee and radar power splits & 9/24 & 6 to 8, 12, 16, 17 & Reading the port number instead of the sign pattern \\
4 & Two magic tees joined; terminated arms & 5/24 & 9 to 11 & Not re-deriving after termination; use \code{sparam.terminate()} \\
2 & Single-stub shunt tuner & 6/24 & 2 to 7 & Two valid answers exist; give both $d$ and $\ell$ pairs \\
6 & Amplifier matching networks and gain & 10/24 & 5 to 8 & Conjugate match at one port only; check $\Gamma_S = \Gamma_{in}^*$ \\
4 & Identify the device from its S-matrix & 4/24 & 13 to 15 & Every one of them is a magic tee, printed without its $1/\sqrt{2}$ or with the arms swapped \\
3 & Test a printed S-matrix; amplitude modulator & 3/24 & 1 to 3 & $RL$ uses $20\log$, because $S_{11}$ is an amplitude ratio \\
7 & Radiation zones with numbers & 5/24 & 1, 2 & $D<\lambda$ makes the two formulas invalid; give the small-antenna boundaries \\
4 & Circular guide modes with numbers & 3/24 & 4, 5 & Bessel roots: $p'_{11}=1.841$ for TE, $p_{01}=2.405$ for TM \\
6 & Negative-resistance oscillator & 2/24 & 9 & $|R_{out}| \ge 1.2R_L$ is a rule, not a suggestion \\
\bottomrule
\end{tabularx}

\vspace{2pt}
{\footnotesize\color{sub} \mk{Every published answer in that volume is machine checked.}
\code{rf.py} solves the lines exactly, \code{verify.py} walks each published stub design
forward and asserts $y_{in}=1+j0$ \textbf{re-deriving the susceptance from the published
length}, \code{sparam.py} re-derives each device matrix from its stated properties,
\code{wg.py} tests every printed field component against both curl equations and the walls,
\code{amp.py} and \code{filt.py} reproduce the deck examples, \code{tubes.py} simulates the
electrons, \code{ant.py} and \code{meas.py} assert every number in Chapters 7 and 8.}

\par\penalty0\vspace{4pt}

%=======================================================================
\T{R.4 Most Repeated PYQs \mk{(if you revise only fifteen questions, revise these)}}

\begin{enumerate}[label=\arabic*., leftmargin=1.6em, itemsep=1.5pt]
  \item \textbf{Check the stability of a transistor with these S-parameters and compare
    maximum power gain for bilateral and unilateral modes.} \yr{22 papers} $\rightarrow$
    $\S$6.4 and Problems 1 to 4.
  \item \textbf{Why is the insertion loss method used for microwave filters? Prototype a
    Butterworth or Chebyshev LPF and implement it in microstrip.} \yr{20 papers}
    $\rightarrow$ $\S$6.1 to $\S$6.3.
  \item \textbf{Differentiate the behaviour of microwave circuits from conventional
    low-frequency circuits.} \yr{18 papers} $\rightarrow$ $\S$1.2.
  \item \textbf{How are microwaves hazardous to humans?} \yr{17 papers} $\rightarrow$
    $\S$7.2.
  \item \textbf{Design a double-stub matching network for this load and prepare the S-matrix
    of the matched network.} \yr{17 papers} $\rightarrow$ $\S$2.7, $\S$2.8, Problems 8 to 24.
  \item \textbf{Derive the field equations and characteristic parameters of a rectangular
    waveguide in TE (or TM) mode.} \yr{14 papers} $\rightarrow$ $\S$4.1.
  \item \textbf{Why are S-parameters used in microwave network analysis? Compare with
    h-parameters.} \yr{12 papers in Ch 3, 10 in Ch 1} $\rightarrow$ $\S$3.1, $\S$1.3.
  \item \textbf{Classify the microwave radiation fields around an antenna, oven or BTS.}
    \yr{11 papers} $\rightarrow$ $\S$7.1 and its two problems.
  \item \textbf{Give detailed protocols and specifications of three (or five) microwave
    communication systems.} \yr{11 papers} $\rightarrow$ $\S$1.5.
  \item \textbf{Draw a neat diagram of a magic tee and derive its S-matrix.}
    \yr{9 papers} $\rightarrow$ $\S$4.5.
  \item \textbf{What is bunching effect? Explain velocity modulation.} \yr{9 papers}
    $\rightarrow$ $\S$5.2. Usually the 2-mark rider on a tube question.
  \item \textbf{Sketch a cross-sectional magnetron with $45^\circ$ (or $60^\circ$,
    $90^\circ$) phase shift between adjacent cavities and explain its working.}
    \yr{9 papers} $\rightarrow$ $\S$5.5.
  \item \textbf{Describe the construction and operation of a multicavity klystron.}
    \yr{7 papers} $\rightarrow$ $\S$5.3.
  \item \textbf{State and justify the properties of the scattering matrix; analyse a
    three-port network.} \yr{6 papers} $\rightarrow$ $\S$3.3, $\S$3.5.
  \item \textbf{Describe how standing waves and microwave power are measured (VSWR meter
    plus bolometry or calorimetry).} \yr{the pairing itself in 72 Ash, 70 Bh and 79 Bh; some
    Chapter 8 question in 22 of 24 papers} $\rightarrow$ $\S$8.3, $\S$8.5, $\S$8.6.
\end{enumerate}

\par\penalty0\vspace{4pt}

%=======================================================================
\T{R.5 One-Day Revision Checklist}

\lead{Morning: the two heavy chapters (about 4 hrs)}
\begin{itemize}
  \item \textbf{Chapter 6} (22.4\,\% of every mark ever set). Work Problem 1 end to end
    without looking: $\Delta$, $K$, $|\Delta|$, verdict, $G_{T,max}$. Then Problem 10's
    ladder, and draw the six named microstrip layouts from memory.
  \item \textbf{Chapter 2} (15.6\,\%). Draw one double-stub construction from a blank chart:
    SWR circle, spacing circle rotated $d$, constant-$g$ arc, both stub lengths off the rim.
    Say the two chart conventions out loud before you start.
  \item If either feels shaky, stop here. These two are a third of the paper.
\end{itemize}

\lead{Afternoon: the reliable marks (about 3 hrs)}
\begin{itemize}
  \item \textbf{Chapter 4}: write the six derivation steps and the characteristic-equation
    table. Then the magic tee matrix, identify it from a printed matrix by its signs, and
    one two-tee junction.
  \item \textbf{Chapter 7}: the zone figure, the hazard chain, HERP / HERO / HERF, SAR and
    the three limits. This is a guaranteed question and it is fast.
  \item \textbf{Chapter 3}: the four tests, the 3-port impossibility theorem, and the
    five-step S-parameter derivation.
  \item \textbf{Chapter 8}: the power decision table, the two calorimeter formulas, and the
    bolometer double bridge. Three answers cover 22 papers.
\end{itemize}

\lead{Evening: consolidation (about 1 hr)}
\begin{itemize}
  \item \textbf{Chapter 5}: velocity modulation, bunching, klystron, magnetron, TWT. All
    descriptive, all the same three-line skeleton: velocity-modulate, bunch, decelerate.
  \item \textbf{Chapter 1}: the band table, the $d/\lambda$ criterion, and the five systems
    with their specifications.
  \item Re-read R.2 straight through. Those are the lines that cannot be rebuilt in the
    hall.
\end{itemize}

\lead{In the hall}
\begin{itemize}
  \item \textbf{Read the marks split first.} A $[2+8]$ is two answers; give the 2-mark
    definition its own line and heading, then the long part.
  \item \textbf{Draw first, write second.} Nearly every question here has one figure that
    earns marks on its own: the chart construction, the guide with $a$ and $b$, the magic
    tee, the zone diagram, the bridge circuit.
  \item \textbf{Say which case you are in} before computing: $K>1$ or not, $a>b$ or not,
    $D>\lambda$ or not, which power band. Most lost marks in this subject are a correct
    method applied to the wrong case.
  \item If a number will not come out, \textbf{write the formula and the substitution
    anyway}. In this subject the method carries most of the mark.
\end{itemize}
